%! AP Statistics — Unit 1, Lesson 1.3 — Slides (Beamer)
\documentclass[aspectratio=169, 11pt]{beamer}
\usepackage{apstats-beamer}

%=============================================================================
\begin{document}

% ─────────────────────────────────────────────────────────────────────────────
% SLIDE 1 — LESSON TITLE
% ─────────────────────────────────────────────────────────────────────────────
{
\setbeamercolor{background canvas}{bg=navy}
\begin{frame}[plain]
  \begin{tikzpicture}[remember picture, overlay]
    \fill[navylight] (current page.north west)
      rectangle ([yshift=-1.35cm]current page.north east);
  \end{tikzpicture}
  \vspace{2.8cm}
  \hspace{0.55cm}%
  \begin{minipage}{0.9\textwidth}
    {\color{goldacc}\bfseries\small LESSON 1.3}\\[0.5cm]
    {\color{white}\bfseries\LARGE Representing a Categorical Variable\\[0.25cm]
      with Tables}\\[1.4cm]
    {\color{skymid}\small AP Statistics~$\cdot$~Unit 1: Exploring One-Variable Data}
  \end{minipage}
\end{frame}
}

% ─────────────────────────────────────────────────────────────────────────────
% SLIDE 2 — WARM-UP
% ─────────────────────────────────────────────────────────────────────────────
{
\setbeamercolor{background canvas}{bg=sky}
\begin{frame}[t, plain]
  \navyheader{Warm-Up}
  \hspace{0.55cm}%
  \begin{minipage}{0.9\textwidth}
    \vspace{0.3em}
    \small

    \begin{enumerate}
      \setlength{\itemsep}{0.7em}\setlength{\topsep}{0em}
      \item Look at the list of 20 students' favorite subjects on your warm-up sheet.
            Can you immediately tell which subject is most popular? Why or why not?
      \item Tally each subject and complete the frequency column.
            Does your total equal 20?
      \item From Lesson 1.2: what \textit{type} of variable is ``favorite subject''?
            Is it meaningful to compute the average?
      \item What question do you think today's lesson will answer?
    \end{enumerate}

    \vspace{0.5cm}
    {\color{navy}\itshape\small
      We'll revisit your frequency table at the end of the lesson to add
      one more column.}
  \end{minipage}
\end{frame}
}

% ─────────────────────────────────────────────────────────────────────────────
% SLIDE 3 — PRIMARY OBJECTIVE
% ─────────────────────────────────────────────────────────────────────────────
{
\setbeamercolor{background canvas}{bg=sky}
\begin{frame}[t, plain]
  \begin{tikzpicture}[remember picture, overlay]
    \fill[navy] (current page.north west)
      rectangle ([yshift=-0.25cm]current page.north east);
  \end{tikzpicture}
  \vspace{0.45cm}\par
  \hspace{0.55cm}%
  \begin{minipage}{0.9\textwidth}

    \sectionlabel{Primary Objective}

    \normalsize
    Students will construct and interpret \textbf{frequency} and
    \textbf{relative frequency} tables for a categorical variable.
    Students will explain why relative frequency --- not raw count --- is
    required to compare groups of different sizes (Big Idea: UNC).

    \vspace{0.45cm}

    \normalsize\textbf{\color{navy}By the end of this lesson, I will be able to:}

    \smallskip
    \begin{itemize}
      \setlength{\itemsep}{0.3em}\setlength{\topsep}{0.2em}
      \item Construct a frequency and relative frequency table from raw categorical data.
      \item Interpret relative frequencies in context using complete sentences.
      \item Explain why relative frequency is necessary for fair comparisons.
    \end{itemize}

  \end{minipage}
\end{frame}
}

% ─────────────────────────────────────────────────────────────────────────────
% SLIDE 4 — HOOK
% ─────────────────────────────────────────────────────────────────────────────
{
\setbeamercolor{background canvas}{bg=hookbg}
\begin{frame}[t, plain]
  \begin{tikzpicture}[remember picture, overlay]
    \fill[goldacc] (current page.north west)
      rectangle ([yshift=-1.35cm]current page.north east);
    \node[anchor=west, text=white, font=\bfseries\large,
          inner xsep=0.55cm, inner ysep=0pt]
      at ([yshift=-0.68cm]current page.north west)
      {Hook --- Which School Has More Walkers?};
  \end{tikzpicture}
  \vspace{1.2cm}\par
  \hspace{0.55cm}%
  \begin{minipage}{0.9\textwidth}
    \vspace{0.2cm}
    \small
    {\color{white}\bfseries
      Two schools survey students about how they get to school.
      Which school has more walkers?}

    \vspace{0.45cm}

    \setlength{\tabcolsep}{10pt}
    \renewcommand{\arraystretch}{1.6}
    \begin{tabular}{@{} >{\bfseries\color{white}}p{0.16\textwidth}
                        >{\centering\arraybackslash}p{0.18\textwidth}
                        >{\centering\arraybackslash}p{0.18\textwidth} @{}}
      \rowcolor{navy}
      {\color{skymid}Mode} &
      {\color{skymid}School A} &
      {\color{skymid}School B} \\
      \rowcolor{white!15!hookbg}
      {\color{white}Bus}  & {\color{white}120} & {\color{white}30} \\
      \rowcolor{white!8!hookbg}
      {\color{white}Car}  & {\color{white}80}  & {\color{white}50} \\
      \rowcolor{white!15!hookbg}
      {\color{white}Walk} & {\color{white}40}  & {\color{white}120} \\
      \rowcolor{navy}
      {\color{skymid}\textbf{Total}} &
      {\color{skymid}\textbf{240}} &
      {\color{skymid}\textbf{200}} \\
    \end{tabular}

    \vspace{0.5cm}
    {\color{white}\small
      Discuss with your partner: Is it fair to compare the ``Walk'' counts directly?}
  \end{minipage}
\end{frame}
}

% ─────────────────────────────────────────────────────────────────────────────
% SLIDE 5 — HOOK REVEAL
% ─────────────────────────────────────────────────────────────────────────────
{
\setbeamercolor{background canvas}{bg=hookbg}
\begin{frame}[t, plain]
  \begin{tikzpicture}[remember picture, overlay]
    \fill[goldacc] (current page.north west)
      rectangle ([yshift=-1.35cm]current page.north east);
    \node[anchor=west, text=white, font=\bfseries\large,
          inner xsep=0.55cm, inner ysep=0pt]
      at ([yshift=-0.68cm]current page.north west)
      {Hook --- The Answer Changes with Relative Frequency};
  \end{tikzpicture}
  \vspace{1.2cm}\par
  \hspace{0.55cm}%
  \begin{minipage}{0.9\textwidth}
    \small

    \setlength{\tabcolsep}{8pt}
    \renewcommand{\arraystretch}{1.5}
    \begin{tabular}{@{} >{\bfseries\color{white}}p{0.14\textwidth}
                        >{\centering\arraybackslash}p{0.12\textwidth}
                        >{\centering\arraybackslash}p{0.16\textwidth}
                        >{\centering\arraybackslash}p{0.12\textwidth}
                        >{\centering\arraybackslash}p{0.16\textwidth} @{}}
      \rowcolor{navy}
      {\color{skymid}Mode} &
      {\color{skymid}Sch.\ A} & {\color{skymid}Rel.\ Freq.\ A} &
      {\color{skymid}Sch.\ B} & {\color{skymid}Rel.\ Freq.\ B} \\
      \rowcolor{white!15!hookbg}
      {\color{white}Bus}  & {\color{white}120} & {\color{white}50\%} &
        {\color{white}30}  & {\color{white}15\%} \\
      \rowcolor{white!8!hookbg}
      {\color{white}Car}  & {\color{white}80}  & {\color{white}33\%} &
        {\color{white}50}  & {\color{white}25\%} \\
      \rowcolor{white!15!hookbg}
      {\color{white}\textbf{Walk}} &
        {\color{white}40}  & {\color{white}\textbf{17\%}} &
        {\color{white}120} & {\color{goldacc}\textbf{60\%}} \\
    \end{tabular}

    \vspace{0.5cm}
    {\color{white}\small
      School A has more walkers by \textit{count} (40 vs.\ 30).\\
      School B has far more walkers by \textit{proportion} (60\% vs.\ 17\%).\\[8pt]
      \textbf{Lesson:} Use relative frequency to make fair comparisons.}
  \end{minipage}
\end{frame}
}

% ─────────────────────────────────────────────────────────────────────────────
% SLIDE 6 — VOCABULARY
% ─────────────────────────────────────────────────────────────────────────────
{
\setbeamercolor{background canvas}{bg=greenbg}
\begin{frame}[t, plain]
  \navyheader{Vocabulary}
  \hspace{0.55cm}%
  \begin{minipage}{0.9\textwidth}
    \small
    \begin{description}[leftmargin=0em, itemsep=0.45em, font=\bfseries\color{navy}]
      \item[Frequency] The count of individuals in a particular category.
      \item[Relative frequency] Frequency $\div$ total; expressed as a proportion
            (decimal) or percent.
      \item[Frequency table] A table listing each category of a categorical variable
            and its count.
      \item[Relative frequency table] A table listing each category and its relative
            frequency.
      \item[Distribution (categorical)] The set of possible values and their associated
            relative frequencies.
    \end{description}

    \vspace{0.4cm}
    {\color{navy}\small\textit{Key check: Relative frequencies must sum to 1 (or 100\%).}}
  \end{minipage}
\end{frame}
}

% ─────────────────────────────────────────────────────────────────────────────
% SLIDE 7 — BUILDING A TABLE
% ─────────────────────────────────────────────────────────────────────────────
{
\setbeamercolor{background canvas}{bg=white}
\begin{frame}[t, plain]
  \navyheader{Building a Relative Frequency Table --- Four Steps}
  \hspace{0.55cm}%
  \begin{minipage}{0.9\textwidth}
    \small

    \begin{enumerate}
      \setlength{\itemsep}{0.55em}
      \item \textbf{List all categories} in the first column.
      \item \textbf{Tally} the raw data and record the \textbf{frequency} (count).
      \item \textbf{Divide} each count by $n$ (the total) to get the
            \textbf{relative frequency}.
      \item \textbf{Check:} all relative frequencies must sum to exactly 1.
    \end{enumerate}

    \vspace{0.5cm}
    \[
      \text{Relative frequency} = \frac{\text{frequency}}{n}
    \]

    \vspace{0.3cm}
    {\color{navy}\small
      Express as a decimal \textit{or} a percent --- both are correct on the AP exam.}
  \end{minipage}
\end{frame}
}

% ─────────────────────────────────────────────────────────────────────────────
% SLIDE 8 — WORKED EXAMPLE
% ─────────────────────────────────────────────────────────────────────────────
{
\setbeamercolor{background canvas}{bg=white}
\begin{frame}[t, plain]
  \navyheader{Worked Example --- Pet Ownership ($n = 40$)}
  \hspace{0.55cm}%
  \begin{minipage}{0.9\textwidth}
    \small

    \renewcommand{\arraystretch}{1.5}
    \begin{tabular}{@{} >{\bfseries}p{0.20\textwidth}
                        >{\centering\arraybackslash}p{0.14\textwidth}
                        >{\centering\arraybackslash}p{0.22\textwidth} @{}}
      \rowcolor{sky}
      \textbf{Pet type} & \textbf{Frequency} & \textbf{Rel.\ Frequency} \\
      \midrule
      Dog   & 12 & $12/40 = 0.30$ \\
      Cat   & 8  & $8/40  = 0.20$ \\
      Fish  & 6  & $6/40  = 0.15$ \\
      Bird  & 4  & $4/40  = 0.10$ \\
      None  & 10 & $10/40 = 0.25$ \\
      \midrule
      \rowcolor{sky}\textbf{Total} & 40 & 1.00 \\
    \end{tabular}

    \vspace{0.5cm}
    {\color{navy}\bfseries\small AP-style interpretation:}\\[4pt]
    \small\textit{``Of the 40 households surveyed, the most common pet type was
    dog, owned by 30\% of households. The least common was bird, owned by 10\%.''}
  \end{minipage}
\end{frame}
}

% ─────────────────────────────────────────────────────────────────────────────
% SLIDE 9 — AP LANGUAGE
% ─────────────────────────────────────────────────────────────────────────────
{
\setbeamercolor{background canvas}{bg=sky}
\begin{frame}[t, plain]
  \navyheader{Describing a Distribution --- AP Language}
  \hspace{0.55cm}%
  \begin{minipage}{0.9\textwidth}
    \small

    \textbf{\color{navy}Every description must include:}
    \begin{itemize}
      \setlength{\itemsep}{0.4em}
      \item The \textbf{total $n$} (``of the 40 households surveyed\ldots'')
      \item The \textbf{variable} in context (``pet type'')
      \item The \textbf{most} and \textbf{least} common categories
      \item The \textbf{relative frequency} for at least the key categories
    \end{itemize}

    \vspace{0.5cm}
    {\color{navy}\bfseries\small Model sentence frame:}\\[6pt]
    \small\textit{``Of the \underline{\hspace{1cm}} [individuals] surveyed, the most
    common [variable] was \underline{\hspace{1.5cm}}, with \underline{\hspace{1cm}}\%
    of responses. The least common was \underline{\hspace{1.5cm}} with
    \underline{\hspace{1cm}}\%.''}

    \vspace{0.4cm}
    {\color{slate}\footnotesize
      On the AP exam, writing ``30\%'' without context earns \textbf{no credit}.
      Always anchor your numbers.}
  \end{minipage}
\end{frame}
}

% ─────────────────────────────────────────────────────────────────────────────
% SLIDE 10 — GROUP ACTIVITY
% ─────────────────────────────────────────────────────────────────────────────
{
\setbeamercolor{background canvas}{bg=redbg}
\begin{frame}[t, plain]
  \begin{tikzpicture}[remember picture, overlay]
    \fill[redacc] (current page.north west)
      rectangle ([yshift=-1.35cm]current page.north east);
    \node[anchor=west, text=white, font=\bfseries\large,
          inner xsep=0.55cm, inner ysep=0pt]
      at ([yshift=-0.68cm]current page.north west)
      {Group Activity --- Favorite Sport Survey};
  \end{tikzpicture}
  \vspace{1.2cm}\par
  \hspace{0.55cm}%
  \begin{minipage}{0.9\textwidth}
    \small

    \textbf{Part 1 (8 min):} Raw data for 30 middle school students.
    \begin{enumerate}
      \setlength{\itemsep}{0.3em}
      \item Tally the data and build a complete frequency and relative frequency table.
      \item Verify: total = 30, sum of relative frequencies = 1.
      \item Write 2--3 sentences describing the distribution.
    \end{enumerate}

    \vspace{0.5cm}
    \textbf{Part 2 (5 min):} A second dataset has $n = 50$.
    \begin{enumerate}
      \setlength{\itemsep}{0.3em}
      \item Complete the relative frequency column for the high school data.
      \item Compare the two groups. Which sport is more popular proportionally?
      \item Can you compare using counts alone? Why or why not?
    \end{enumerate}
  \end{minipage}
\end{frame}
}

% ─────────────────────────────────────────────────────────────────────────────
% SLIDE 11 — EXIT TICKET
% ─────────────────────────────────────────────────────────────────────────────
{
\setbeamercolor{background canvas}{bg=sky}
\begin{frame}[t, plain]
  \navyheader{Exit Ticket --- 5 Minutes}
  \hspace{0.55cm}%
  \begin{minipage}{0.9\textwidth}
    \small

    A survey of \textbf{50 students} asks: ``Where do you prefer to study?''

    \vspace{0.35cm}
    \renewcommand{\arraystretch}{1.5}
    \begin{tabular}{@{} >{\bfseries}p{0.22\textwidth}
                        >{\centering\arraybackslash}p{0.18\textwidth}
                        >{\centering\arraybackslash}p{0.22\textwidth} @{}}
      \rowcolor{navy}
      {\color{white}Location} & {\color{white}Count} & {\color{white}Rel.\ Freq.} \\
      \midrule
      Library  & 18 & ? \\
      Home     & 22 & ? \\
      Caf\'{e} & 6  & ? \\
      Other    & 4  & ? \\
      \rowcolor{sky}\textbf{Total} & 50 & ? \\
    \end{tabular}

    \vspace{0.5cm}
    \begin{enumerate}
      \setlength{\itemsep}{0.4em}
      \item Complete the relative frequency column.
      \item Write one sentence describing the most common location \textit{in context}.
      \item What percentage of students study somewhere other than home?
    \end{enumerate}
  \end{minipage}
\end{frame}
}

% ─────────────────────────────────────────────────────────────────────────────
% SLIDE 12 — BIG IDEAS / PREVIEW
% ─────────────────────────────────────────────────────────────────────────────
{
\setbeamercolor{background canvas}{bg=navy}
\begin{frame}[t, plain]
  \begin{tikzpicture}[remember picture, overlay]
    \fill[navylight] (current page.north west)
      rectangle ([yshift=-1.35cm]current page.north east);
    \node[anchor=west, text=white, font=\bfseries\large,
          inner xsep=0.55cm, inner ysep=0pt]
      at ([yshift=-0.68cm]current page.north west)
      {Big Ideas \& Preview};
  \end{tikzpicture}
  \vspace{1.2cm}\par
  \hspace{0.55cm}%
  \begin{minipage}{0.9\textwidth}
    \small

    {\color{goldacc}\bfseries Today's core takeaway:}\\[6pt]
    {\color{white}
    Use \textbf{relative frequency} when comparing groups of different sizes.
    Counts alone are misleading.}

    \vspace{0.55cm}
    {\color{skymid}\bfseries Connections:}
    \begin{itemize}
      \setlength{\itemsep}{0.35em}
      \item {\color{white}
            \textbf{Lesson 1.4:} The same distributions shown as bar charts and
            pie charts --- tables vs.\ graphs, trade-offs.}
      \item {\color{white}
            \textbf{Unit 2:} Two-way tables extend today's one-way tables to
            two categorical variables.}
      \item {\color{white}
            \textbf{Unit 4:} Relative frequency $=$ experimental probability ---
            the arithmetic is identical.}
    \end{itemize}

    \vspace{0.55cm}
    {\color{skymid}\small\itshape
      Complete your homework (Lesson 1.3 HW) before next class.}
  \end{minipage}
\end{frame}
}

\end{document}
